\documentclass[11pt,reqno]{amsart}

%=======================================================================
%  Structural and Computational Results on the Finite-Support
%  Highway Conjecture for Langton's Ant
%  Author: Atharva Jillhewar
%  Compiles with pdfLaTeX on Overleaf (standard packages only).
%=======================================================================

\usepackage[utf8]{inputenc}
\usepackage[T1]{fontenc}
\usepackage{lmodern}      % Type 1 vector fonts: keeps the PDF text layer searchable/copyable
\usepackage{microtype}
\usepackage{amsmath,amssymb,amsthm,mathtools}
\usepackage{enumitem}
\usepackage{booktabs}
\usepackage{tikz}
\usetikzlibrary{arrows.meta}
\usepackage[margin=1in]{geometry}
\usepackage[colorlinks=true,linkcolor=blue!55!black,citecolor=teal!60!black,urlcolor=purple!60!black]{hyperref}

%------------------------------ theorem environments
\theoremstyle{plain}
\newtheorem{theorem}{Theorem}[section]
\newtheorem{lemma}[theorem]{Lemma}
\newtheorem{proposition}[theorem]{Proposition}
\newtheorem{corollary}[theorem]{Corollary}
\theoremstyle{definition}
\newtheorem{definition}[theorem]{Definition}
\newtheorem{remark}[theorem]{Remark}
\newtheorem{example}[theorem]{Example}
\newtheorem{conjecture}[theorem]{Conjecture}
\newtheorem{problem}[theorem]{Problem}

%------------------------------ macros
\newcommand{\Z}{\mathbb{Z}}
\newcommand{\F}{\mathbb{F}}
\newcommand{\N}{\mathbb{N}}
\newcommand{\bB}{\mathbf{B}}
\newcommand{\ant}{\mathsf{a}}
\newcommand{\che}{\chi}                 % checker sign
\newcommand{\drift}{d}
\newcommand{\growth}{g}
\newcommand{\turnR}{\mathrm{R}}
\newcommand{\turnL}{\mathrm{L}}
\newcommand{\bd}{\partial}
\DeclareMathOperator{\rank}{rank}
\DeclareMathOperator{\supp}{supp}
\DeclareMathOperator{\dist}{dist}
\newcommand{\half}{\tfrac12}
\newcommand{\PN}[1]{\textbf{P#1}}

\title[Finite-support Langton ant]{Structural and Computational Results on the Finite-Support Highway Conjecture for Langton's Ant}
\author{Atharva Jillhewar}
\date{21 July 2026}
\hypersetup{
  pdftitle={Structural and Computational Results on the Finite-Support Highway Conjecture for Langton's Ant},
  pdfauthor={Atharva Jillhewar}
}

\begin{document}

\begin{abstract}
Langton's ant is a deterministic reversible lattice walker on $\Z^2$ whose
long-term behaviour, for every finite initial colouring, is empirically the
formation of a period-$104$ ``highway''. Whether this holds for \emph{every}
finite initial black support is a well-known open problem. This paper collects,
in a uniform and rigorous framework, a body of structural and computational
results on this conjecture. We prove: an exact, decidable
necessary-and-sufficient criterion for an infinite periodic turn word to be the
trace of some finite-support colouring (Theorem~\ref{thm:criterion}); a
heading/black-count mod-four invariant and a reduction of the two-dimensional
orbit to a one-dimensional black-count walk; an exact conjugacy of the ant with
edge-toggling on a Tait (medial) graph, together with a cycle-rank surgery
identity; and two independent proofs that \emph{no finite-support clean
translator with zero net growth exists}---one homological (the
\emph{even-winding theorem}) and one via a \emph{signed mod-four wake-residue}
identity, which further yields the strand-density bound $\growth\ge
2\|\drift\|_\infty$ for every finite-support repeating trace of nonzero drift.
Consequently every
finite-support periodic highway has net black growth a positive multiple of
four.  For diagonal drift we prove extremal-line rigidity and exclude transverse
width two at every period.  A computer-assisted crossing-sequence theorem then
excludes width four at every period: two independent exact enumerators produce the
same twelve-edge blank-column graph, and Lean checks that every edge lowers rank.
Thus a diagonal periodic highway has transverse width at least six. On the
computational side we give an exact finite black seed for the
standard highway and exclude by a rank-audited exhaustive search---using no
consequence of the residue identity, and
applying the exact criterion to every positive-growth nonzero-drift leaf---every
finite-support periodic highway of nonzero drift and period at most $48$. We
isolate two logically separate remaining problems: a universal
entrance/bounded-retreat problem and the all-period classification of
positive-growth periodic traces. We describe a Lean~4 formalization of selected
algebraic kernels. No
claim of a proof or disproof of the highway conjecture is made.
\end{abstract}

\maketitle

\newpage
\tableofcontents
\newpage

%=======================================================================
\section{Introduction}
%=======================================================================

\subsection{The automaton and the conjecture}
Langton's ant moves on the square lattice $\Z^2$, each cell coloured
\emph{white} ($0$) or \emph{black} ($1$). A pointer (the \emph{ant}) occupies a
cell and faces one of the four cardinal directions. One step of the dynamics is:
\begin{quote}
if the current cell is white, turn right ($90^\circ$ clockwise); if it is black,
turn left; then flip the colour of the current cell and move one unit forward in
the new heading.
\end{quote}
We always assume the \emph{finite-support} condition: the set of initially black
cells is finite. Translating and rotating, we normalise the initial ant to be at
the origin facing north.

Iterating from the all-white configuration produces, after a transient of
$9977$ steps, a self-similar diagonal structure---the \emph{highway}---that
repeats with least period $104$ and displaces the ant by $(\pm2,\pm2)$ per
period. Extensive experiment suggests the same eventual behaviour for every
finite initial colouring.

\begin{conjecture}[Finite-support highway conjecture]\label{conj:main}
For every finite initial black set $B_0\subset\Z^2$ and every initial pose, there
is a time after which the ant's future turn trace and path agree with the standard
period-$104$ highway up to lattice translation, quarter-turn rotation, and cyclic
phase shift. Finite cells outside the future path may remain as stationary debris;
equality of the entire colouring is not asserted.
\end{conjecture}

Conjecture~\ref{conj:main} is stronger than \emph{unboundedness} of the
trajectory (that the ant visits infinitely many cells), which is a theorem of
Bunimovich and Troubetzkoy \cite{BunimovichTroubetzkoy1992}. The ant is
computationally universal, and associated prediction problems are undecidable, on
suitable infinite backgrounds \cite{GajardoMoreiraGoles2002}; we draw no inference
from this about the finite-support conjecture beyond the evident point that no
finite cutoff exhausts an infinite range. Gajardo \cite{Gajardo2014symbolic} gives
an algebraic characterisation of the forbidden factors of the ant's trace subshift.
Recent work on \emph{generalised} ants exhibits rules with infinitely many highways
\cite{GajardoLutfallaRao2024} and rules with non-highway emergent behaviour
\cite{Lutfalla2025sideways}; quantitative confinement bounds for the classical ant
are studied in \cite{Etse2026confined}. Conjecture~\ref{conj:main} remains open.

\subsection{What is proved here, and what is not}
This paper is deliberately explicit about status. We prove a number of exact
structural theorems and finite computational theorems; \emph{none} of them
resolves Conjecture~\ref{conj:main}. Section~\ref{sec:gap} states the two
logically separate unresolved problems. Every statement below is labelled as an
exact theorem, a finite computational theorem, a conjecture, or an explicitly
open problem. The main contributions are:

\begin{enumerate}[leftmargin=2.2em,label=(\arabic*)]
\item An exact periodic-trace criterion (Theorem~\ref{thm:criterion}) turning
``is $w^\omega$ a finite-support highway?'' into a finite computation, with an
explicit seed.
\item The heading/black-count invariant and the black-count path reduction
(Section~\ref{sec:invariants}).
\item The Tait-graph / boundary-loop conjugacy and the cycle-rank surgery
identity (Section~\ref{sec:tait}).
\item The even-winding theorem: no zero-growth clean translator
(Theorem~\ref{thm:evenwinding}), hence growth $\in 4\Z_{>0}$ for every highway.
\item The signed mod-four wake-residue theorem
(Theorem~\ref{thm:residue}) and the strand-density bound $\growth\ge
2\max(|a|,|b|)$.
\item Residue-free, rank-audited exclusion of all finite-support periodic
highways of period $\le 48$ (Section~\ref{sec:computation}).
\item A two-endpoint collision-chain parity theorem relevant to the entrance
problem (Section~\ref{sec:interlace}).
\item Extremal-line rigidity and all-period exclusions of diagonal transverse
widths two and four; the width-four result has an explicit computer-assisted local
classification boundary (Section~\ref{sec:transverse}).
\item A Lean~4 formalization of selected algebraic kernels (Appendix~\ref{app:lean}).
\end{enumerate}

The presentation is organised so that failed and merely heuristic arguments are
recorded honestly (Section~\ref{sec:gap}), because they delimit the true
difficulty of the problem.

%=======================================================================
\section{Conventions and elementary invariants}\label{sec:prelim}
%=======================================================================

\subsection{State, step, trace}
A \emph{state} is a triple $s=(B,p,q)$ with $B\subset\Z^2$ finite (the black
set), $p\in\Z^2$ (the ant position) and $q\in\Z/4\Z$ the heading, numbered
$0=N,1=E,2=S,3=W$. Let $u_0=(0,1),u_1=(1,0),u_2=(0,-1),u_3=(-1,0)$ be the unit
heading vectors. The \emph{turn at $s$} is $\turnR$ (right) if $p\notin B$ and
$\turnL$ (left) if $p\in B$; write $\varepsilon=+1$ for $\turnR$ and
$\varepsilon=-1$ for $\turnL$. One step is
\[
s\mapsto\bigl(B\triangle\{p\},\; p+u_{q+\varepsilon},\; q+\varepsilon\bigr),
\]
where $\triangle$ is symmetric difference. The \emph{trace} of an orbit is the
infinite word $w\in\{\turnR,\turnL\}^{\N}$ of turns encountered.

\begin{definition}[Periodic terminology]\label{def:periodic}
Let $w=w_0\cdots w_{N-1}$ be a nonempty turn word and follow its induced lattice
path from a specified initial heading. The word is \emph{heading-resetting} if the
heading after $N$ turns equals the initial heading. In that case its \emph{drift} is
$p_N-p_0$, and its \emph{growth} is
$\#\{i:w_i=\turnR\}-\#\{i:w_i=\turnL\}$. An orbit has an \emph{eventually periodic
turn trace} if, from some time $t_0$, its turns are $w^\omega$. A finite-support
tail of this form with $w$ heading-resetting is called a \emph{repeating trace}; if
its drift is nonzero, it is a \emph{periodic highway}. Its \emph{least trace period}
is the smallest possible $N$. Highways are compared up to lattice translation,
quarter-turn rotation, and cyclic phase shift of $w$.
\end{definition}

The dynamics is reversible: the pre-image of $(B,p,q)$ is obtained by stepping
backward one cell, recolouring, and undoing the turn; hence orbits never merge.

\subsection{HV partition and alternating visits}

\begin{proposition}[\PN{1}: HV/checkerboard partition]\label{prop:hv}
Along any orbit the heading parity $q_t\bmod 2$ and the checkerboard parity
$(p_{t,x}+p_{t,y})\bmod 2$ each change at every step, so their difference is a
constant of the motion. Consequently each cell $z$ is entered by the ant always
horizontally or always vertically; the admissible directed lattice edges have a
fixed orientation, and each cell has exactly two incoming and two outgoing
directed edges.
\end{proposition}

\begin{proof}
Each step turns by $\pm1$ quarter turn, so $q_{t+1}-q_t$ is odd and $q_t\bmod2$
alternates. The displacement $u_{q_{t+1}}$ has odd $L^1$-length, so
$(p_{t,x}+p_{t,y})\bmod 2$ also alternates. Their sum is therefore invariant.
Fix a cell $z$; whenever $p_t=z$ the value of $(p_{t,x}+p_{t,y})\bmod2$ is fixed,
hence so is $q_t\bmod2$: the arrival is always via a horizontal edge or always
via a vertical edge. A vertical-arrival cell has its two vertical edges incoming
and horizontal edges outgoing, and vice versa.
\end{proof}

Call a cell \emph{V-type} if the ant always arrives vertically and \emph{H-type}
otherwise. We write $\che(x,y)=(-1)^{x+y}$ for the checkerboard sign; V/H type is
determined by $\che$ together with the fixed invariant of
Proposition~\ref{prop:hv}.

\begin{proposition}[\PN{2}: alternating visit rule]\label{prop:alt}
At any fixed cell the turns encountered alternate strictly:
$\turnR,\turnL,\turnR,\turnL,\dots$ or $\turnL,\turnR,\turnL,\dots$. Since a cell
is white iff it will next cause an $\turnR$, and only finitely many cells are
initially black, all but finitely many cells have $\turnR$ as their first
encountered turn.
\end{proposition}

\begin{proof}
Each visit flips the colour, so the colour seen at a cell alternates; the turn
is a function of colour, hence alternates too. A cell that is white at its first
visit begins with $\turnR$; only cells in $B_0$ can begin with $\turnL$.
\end{proof}

\begin{theorem}[Unboundedness \cite{BunimovichTroubetzkoy1992}]\label{thm:unbdd}
Every finite-support orbit visits infinitely many cells.
\end{theorem}

\begin{proof}[Proof sketch]
If the visited set were finite, determinism and reversibility would make the
orbit periodic in the full state. Take a visited cell $z$ maximising $x+y$; its
north and east neighbours are unvisited. If $z$ is entered horizontally, it can
only be entered from the west. During a full state period its colour returns, so
its alternating visits include both a white/$\turnR$ and a black/$\turnL$ visit;
the corresponding departures include the unvisited north neighbour. If $z$ is
entered vertically, it can only be entered from the south, and the two colours
force departures including the unvisited east neighbour. Either case is a
contradiction.
\end{proof}

\subsection{The heading/black-count invariant and path reduction}\label{sec:invariants}

Let $B_t,p_t,q_t$ be the state before step $t$, $b_t=|B_t|$, and
$r_t=b_{t+1}-b_t\in\{\pm1\}$.

\begin{proposition}[\PN{4}, \PN{9}: heading/count invariant and path reduction]
\label{prop:pathreduction}
For all $t$,
\begin{equation}\label{eq:P4}
q_t-b_t\equiv q_0-b_0\pmod 4,
\end{equation}
and the two-dimensional path is a function of the one-dimensional black-count
walk:
\begin{equation}\label{eq:P9}
p_{t+1}-p_t=u_{\,q_0-b_0+b_{t+1}}.
\end{equation}
In particular a heading-resetting period changes $b$ by a multiple of $4$.
\end{proposition}

\begin{proof}
A white/$\turnR$ step has $r_t=+1$ and $q_{t+1}-q_t=+1$; a black/$\turnL$ step
has both differences $-1$. Thus $q_t-b_t$ is invariant, giving \eqref{eq:P4}.
The move occurs after the turn, so $q_{t+1}=q_0-b_0+b_{t+1}\pmod4$ and
\eqref{eq:P9} follows.
\end{proof}

Equation \eqref{eq:P9} does not say every nonnegative nearest-neighbour walk is
colour-realisable: repeated visits still must request alternating turns
(Proposition~\ref{prop:alt}). Differentiating the centred first moment
$J_t=\sum_{z\in B_t}(z-p_t)$ gives the exact identity
\begin{equation}\label{eq:P10}
J_{t+1}-J_t=-\,b_{t+1}\,u_{\,q_0-b_0+b_{t+1}},
\end{equation}
which is invariant under simultaneous translation of the black set and the ant
(\PN{10}); it will quantify the wake of a growing highway.

%=======================================================================
\section{The exact periodic-trace criterion}\label{sec:criterion}
%=======================================================================

The following theorem converts the search for a nonstandard highway into a
finite, exact decision problem. Fix an initial heading $q_0$ and a proposed
finite turn word $w=w_0\cdots w_{N-1}$. Compute the associated path
$p_0,p_1,\dots,p_N$ purely from $w$ (Proposition~\ref{prop:hv} makes each step
well defined). Suppose the heading is restored ($q_N=q_0$) and the displacement
\[
\drift=p_N-p_0\neq 0 .
\]
Define an equivalence on phases by $i\sim j\iff p_i-p_j\in\Z\drift$. For each
class $I$ pick a representative $b_I$ and integers $a_i$ with $p_i=b_I+a_i\drift$
($i\in I$). For $n\ge\min_{i\in I}a_i$, let $S_{I,n}$ be the word formed by all
$w_i$ ($i\in I$) with $a_i\le n$, ordered by the pair $(n-a_i,\,i)$.

\begin{theorem}[\PN{3}: finite-support periodic-trace criterion]\label{thm:criterion}
The infinite word $w^\omega$ is the trace of Langton's ant from some
finite-support colouring if and only if, for every class $I$:
\begin{enumerate}[label=\textup{(\roman*)}]
\item every $S_{I,n}$ with $\min_I a_i\le n\le\max_I a_i$ alternates strictly
between $\turnR$ and $\turnL$; and
\item the stabilised word $S_{I,\max_I a_i}$ begins with $\turnR$.
\end{enumerate}
When these hold, an explicit finite black seed is
\[
B=\bigl\{\,b_I+n\drift:\ \textstyle\min_I a_i\le n<\max_I a_i\ \text{and}\
S_{I,n}\ \text{begins with}\ \turnL\,\bigr\},
\]
the union over all classes $I$; the two bounds make $B$ finite and $S_{I,n}$
well defined.
\end{theorem}

\begin{proof}
$S_{I,n}$ is precisely the chronological turn sequence at the physical cell
$b_I+n\drift$ over the infinite trajectory $w^\omega$: the occurrence of phase
$i\in I$ in period $k\ge0$ is at $p_i+k\drift=b_I+(a_i+k)\drift$, which equals
$b_I+n\drift$ exactly when $k=n-a_i\ge0$, and chronological order is the
lexicographic order of $(n-a_i,i)$.

\emph{Necessity.} Each visit flips the cell, so its turn sequence alternates:
(i) is necessary. Once $n\ge\max_I a_i$ every phase of $I$ is present and the
ordering no longer depends on $n$; the sequence has stabilised, and infinitely
many cells share this first turn. By Proposition~\ref{prop:alt} only finitely
many cells may start with $\turnL$, so the stable first turn is $\turnR$: (ii) is
necessary.

\emph{Sufficiency.} Colour exactly $B$ black. At each boundary cell the chosen
colour reproduces the first symbol of $S_{I,n}$, and strict alternation supplies
every later symbol because the cell flips at each visit. At every stabilised cell
(ii) makes the first encounter $\turnR$, matching its white initial colour.
Induction on chronological time shows the ant follows $w^\omega$ forever; the
position is displaced by $\drift\neq0$ every $N$ steps, so this is a highway.
\end{proof}

\begin{remark}\label{rem:onesort}
There is an equivalent one-sequence test: in each class sort all entries by
decreasing $a_i$, breaking ties by increasing $i$; every $S_{I,n}$ is a suffix of
this stabilised word obtained by deleting complete high-level blocks, so it
suffices to check that the single stabilised word alternates and begins with
$\turnR$. To index the classes of $\Z^2/\Z\drift$ exactly, write
$g=\gcd(a,b)$ and $\drift=g(a',b')$ with $\gcd(a',b')=1$, and choose integers $r,s$
with $ra'+sb'=1$. Then
\[
\bigl(b'x-a'y,\ (rx+sy)\bmod g\bigr)
\]
is a complete invariant of the coset of $(x,y)$. Indeed, the matrix
$\left(\begin{smallmatrix}b'&-a'\\ r&s\end{smallmatrix}\right)$ is unimodular, and
in its coordinates translation by $\drift$ becomes translation by $(0,g)$. This is
the checker used in all computations of Section~\ref{sec:computation}.
\end{remark}

%=======================================================================
\section{The Tait-graph / boundary-loop conjugacy}\label{sec:tait}
%=======================================================================

Normalise the HV partition (Proposition~\ref{prop:hv}) so that even cells receive
vertical arrivals. In the \emph{frozen} all-white routing---the permutation on
directed states that ignores colour flips---the admissible directed arcs form
clockwise directed four-cycles around one checkerboard class of unit plaquettes.
Take those plaquettes, labelled by their lower-left corners, as the vertices of a
rotated square lattice. Each ant cell $z=(x,y)$ is the crossing of exactly two
white four-cycles and corresponds to the Tait edge
\begin{equation}\label{eq:tait}
e_z=\begin{cases}
\{(x,y-1),(x-1,y)\}, & x+y \text{ even},\\
\{(x-1,y-1),(x,y)\}, & x+y \text{ odd}.
\end{cases}
\end{equation}
Let $H_B=\{e_z:z\in B\}$ be the finite edge subgraph of black cells.

\begin{theorem}[\PN{5}: boundary-loop conjugacy]\label{thm:conjugacy}
White and black are the two smoothings of the corresponding directed crossing.
Consequently:
\begin{enumerate}[label=\textup{(\alph*)}]
\item the frozen routing loops are exactly the boundary components of a regular
neighbourhood of $H_B$, together with the untouched four-cycles;
\item the ant traverses its current boundary loop and then toggles $e_z$; and
\item every step either merges two frozen loops into one or splits one loop into
two.
\end{enumerate}
This is an exact conjugacy, not a heuristic.
\end{theorem}

\begin{proof}
Frozen (colour-ignoring) routing is the permutation of directed edges sending each
edge to the next one the ant would follow under a fixed local turn rule. At each
cell $z$, right and left turns are precisely the two smoothings of the associated
$4$-valent medial crossing, and the Tait edge $e_z$ records which diagonal pair of
plaquette corners the crossing joins.

Freeze a colouring $B$. Away from $H_B$ every crossing has the white smoothing.
Those smoothings assemble into the boundary circuits of a regular neighbourhood
$\mathcal N$ of $H_B$: occupied edges are thickened to bands and incident Tait
vertices to disks, while untouched white plaquettes contribute trivial four-cycles.
Thus the frozen loops are the components of $\partial\mathcal N$ together with those
trivial cycles. The live ant advances on its current boundary loop and then toggles
$z$, which adds or deletes the single band $e_z$. This surgery reconnects two local
boundary arcs: it splits a component if both arcs belonged to the same loop and
merges them if they belonged to different loops. These are the only cases.
\end{proof}

Let $E=|B|$ be the number of black (occupied) Tait edges, $V$ the number of
incident Tait vertices, $k$ the number of nonempty connected components, and
$\beta=E-V+k$ the cycle rank.

\begin{theorem}[\PN{6}: cycle-rank surgery identity]\label{thm:cyclerank}
A regular neighbourhood of $H_B$ has $k+\beta=E-V+2k$ boundary components.
Relative to the $V$ separate all-white plaquette loops, the change in loop count
is
\begin{equation}\label{eq:cyclerank}
\Delta C=E-2V+2k,\qquad E+\Delta C=2\beta .
\end{equation}
Adding a bridge merges two loops and preserves $\beta$; adding a cycle edge
splits a loop and raises $\beta$ by one. Removing a bridge splits a loop and
preserves $\beta$; removing a nonbridge merges two loops and lowers $\beta$ by
one.
\end{theorem}

\begin{proof}
A regular neighbourhood of a plane graph with $V$ vertices, $E$ edges and $k$
components is a compact planar surface with $k$ genus-zero pieces and Euler
characteristic $V-E$. If it has $b$ boundary circles, then $V-E=2k-b$, so
$b=E-V+2k=k+\beta$, where $\beta=E-V+k$ is the first Betti number.
Formula~\eqref{eq:cyclerank} is the difference from the $V$ trivial plaquette
loops. The surgery cases are the standard effect of adding/deleting a bridge
(does not change $\beta$; changes the loop count by $\mp1$ through
merge/split) versus a cycle edge (changes $\beta$ by $\pm1$; opposite
merge/split), read off from Theorem~\ref{thm:conjugacy}(c).
\end{proof}

\begin{proposition}[\PN{7}: eventual record edges are bridges]\label{prop:record}
Once a new bounding-box record lies strictly beyond the bounding box of the finite
initial support, its first visit is white, and a new east, north, west, or south
record turns south, east, north, or west, respectively. Its Tait edge has a
previously untouched endpoint, so it is a bridge. Hence sufficiently advanced
record creation cannot increase $\beta$: any unbounded production of cycle rank
occurs strictly inside already-explored territory.
\end{proposition}

\begin{proposition}[No monotone affine invariant]\label{rem:nomonotone}
No nonconstant affine combination of $(E,V,k,\beta)$ is monotone along the blank
orbit. In particular, each of $E$, $V$, $k$, and $\beta$ separately fails to be
monotone. Thus no affine combination of these four graph statistics can serve as
a monotone or Lyapunov potential along even the blank orbit.
\end{proposition}

\begin{proof}
Since $\beta=E-V+k$ identically, $aE+bV+ck+d\beta+e=(a+d)E+(b-d)V+(c+d)k+e$: the
affine combinations of $(E,V,k,\beta)$ are exactly those of $(E,V,k)$, nonconstant
precisely when $\lambda:=(a+d,\,b-d,\,c+d)\neq0$. One step flips one cell, hence
inserts or deletes exactly one Tait edge $e_z=\{u,v\}$; insertion has exactly four
cases, and deletion negates each:
\begin{center}
\begin{tabular}{@{}lcc@{}}
\toprule
insertion of $e_z=\{u,v\}$ & $(\Delta E,\Delta V,\Delta k)$ & $\Delta\beta$\\
\midrule
$u,v$ both new & $(1,2,1)$ & $0$\\
exactly one of $u,v$ new & $(1,1,0)$ & $0$\\
$u,v$ both old, distinct components & $(1,0,-1)$ & $0$\\
$u,v$ both old, same component & $(1,0,0)$ & $1$\\
\bottomrule
\end{tabular}
\end{center}
So every increment lies in
$S_0=\pm\{(1,2,1),(1,1,0),(1,0,-1),(1,0,0)\}$, and $\beta$ rises exactly when the
inserted edge closes a cycle. All eight members of $S_0$ occur on the blank orbit
within $12{,}000$ steps. $S_0$ is centrally symmetric by construction and contains
the independent $(1,1,0),(1,0,-1),(1,0,0)$, so it spans $\mathbb{R}^3$. Given
$\lambda\neq0$, spanning yields $v\in S_0$ with $\lambda\cdot v>0$ (replace $v$ by
$-v$ if needed), while $\lambda\cdot(-v)<0$ and $-v$ is realised too: $\lambda$
rises at one step and falls at another. Equivalently $0$ is interior to
$\operatorname{conv}S_0$. See \texttt{verify\_nomonotone.py}.
\end{proof}

%=======================================================================
\section{Charges and clean-translator restrictions}\label{sec:charges}
%=======================================================================

\begin{definition}\label{def:translator}
A \emph{clean finite translator with nonzero drift} is a finite state such that
for some $N>0$ and $\drift\in\Z^2\setminus\{0\}$,
\[
B_N=B_0+\drift,\qquad p_N=p_0+\drift,\qquad q_N=q_0.
\]
``Clean'' means the complete black set translates: no stationary debris, no
newly emitted wake. This is \emph{not} the ordinary highway, which grows a wake.
\end{definition}

\subsection{Additive \texorpdfstring{$\F_2$}{F2} charges and axis alignment}
Work in the Laurent ring $\F_2[X^{\pm1},Y^{\pm1}]$. Let $V_t=1$ if $q_t$ is
vertical, $H_t=1-V_t$, and
\[
A_t=\sum_{(x,y)\in B_t}X^x+V_tX^{x_t},\qquad
D_t=\sum_{(x,y)\in B_t}Y^y+H_tY^{y_t}.
\]

\begin{proposition}[\PN{11}, Step 4]\label{prop:laurent}
$A_t$ and $D_t$ are conserved exactly. Hence a clean translator with
$\drift=(a,b)$ satisfies $A=X^aA$ and $D=Y^bD$ in the integral domain
$\F_2[X^{\pm1},Y^{\pm1}]$; since $A(1)=|B|+V$ and $D(1)=|B|+H$ cannot both vanish
($V+H=1$), $a$ and $b$ cannot both be nonzero. Every clean translator is
axis-aligned.
\end{proposition}

\begin{proof}
For $A_t$: if the incoming heading is vertical the cell toggle at column $x_t$
cancels the disappearing vertical correction; if horizontal, the post-turn
motion is vertical so $x_{t+1}=x_t$ and the toggle cancels the newly appearing
correction. Conservation and translation covariance give $A=X^aA$. In an
integral domain $a\neq0$ forces $A=0$, similarly $b\neq0$ forces $D=0$; but
$A(1)+D(1)=2|B|+1=1\neq0$ in $\F_2$. The row statement is the rotated argument.
\end{proof}

\begin{proposition}[additive-charge cocycle]\label{prop:cocycle}
Let $A$ be an abelian group and $w:\Z^2\to A$. If there exist heading corrections
$F_q:\Z^2\to A$ such that $\sum_{z\in B}w(z)+F_q(p)$ is conserved under
\begin{equation}\label{eq:pot}
F_{q+1}(z+u_{q+1})-F_q(z)=-w(z)\quad(\turnR),\qquad
F_{q-1}(z+u_{q-1})-F_q(z)=+w(z)\quad(\turnL),
\end{equation}
then $w$ satisfies the four-corner equation
\begin{equation}\label{eq:cocycle}
w(x{+}1,y{+}1)+w(x{+}1,y)+w(x,y{+}1)+w(x,y)=0
\end{equation}
and $4w=0$. Conversely, for $A=\Z/4\Z$, the four-corner equation yields such
heading corrections. Over $\F_2$ its solutions are exactly
$w(x,y)=\alpha(x)+\beta(y)$.
\end{proposition}

\begin{proof}
For necessity, eliminate the four $F_q$ from \eqref{eq:pot} around a unit
plaquette to obtain \eqref{eq:cocycle}. At one cell $z$, the
$\turnR$-at-$S$/$\turnL$-at-$N$ relations give $F_S(z)=F_N(z)+2w(z)$, while the
$\turnL$-at-$S$/$\turnR$-at-$N$ relations give $F_S(z)=F_N(z)-2w(z)$; hence
$4w(z)=0$.

For sufficiency over $\Z/4\Z$, compose pairs of equations in \eqref{eq:pot} and
define $F_N$ by the diagonal increments
\begin{align*}
F_N(x{+}1,y{+}1)-F_N(x,y)&=-w(x,y)+w(x{+}1,y),\\
F_N(x{+}1,y{-}1)-F_N(x,y)&=w(x,y{-}1)-w(x{+}1,y{-}1).
\end{align*}
The circulation around each elementary diamond is a sum of two four-corner
expressions plus a multiple of $4w$, so it vanishes. Thus these increments
integrate consistently on both parity sublattices. Set $F_S=F_N-2w$ and define
$F_E,F_W$ from $F_N$ using \eqref{eq:pot}; substitution verifies all eight local
relations, with the only ambiguity equal to $4w=0$. Finally, over $\F_2$, taking
first differences of \eqref{eq:cocycle} shows the horizontal difference depends
only on $x$ (equivalently the vertical difference only on $y$), and integration
gives exactly $w(x,y)=\alpha(x)+\beta(y)$.
\end{proof}

\subsection{The strand decomposition}
Let $\growth=\#\turnR-\#\turnL$ be the net growth of one period of a repeating
word. Restoring the heading forces $\growth\equiv0\pmod4$
(Proposition~\ref{prop:pathreduction}), and $\growth\ge0$ because $\growth$
equals the number of \emph{odd} stabilised translation classes
(Theorem~\ref{thm:criterion}).

\begin{theorem}[\PN{12}: canonical endpoint decomposition]\label{thm:decomp}
Every repeated trace has a canonical decomposition of its one-period signed
toggle pattern
$D(z)=\#\{\turnR\text{ at }z\}-\#\{\turnL\text{ at }z\}$:
\begin{equation}\label{eq:decomp}
D=S+(T_\drift-1)A,\qquad S=\sum_{i=1}^{\growth}\delta_{c_i},
\end{equation}
where $(T_\drift A)(z)=A(z-\drift)$, $A$ is a finite moving widget, and the
$c_i$ are the bases of the $\growth$ growing strands (one $+1$ base per odd
stabilised class).
\end{theorem}

\begin{proof}
$D$ is finitely supported. Partition its support into translation classes under
$z\mapsto z+\drift$. In one class choose a base $c$ and put
$f(n)=D(c+n\drift)$, a finitely supported function on $\Z$. Let
$s=\sum_n f(n)$ and $h=f-s\delta_0$, so $\sum_nh(n)=0$. With
$(Ta)(n)=a(n-1)$, define $a(n)=\sum_{m>n}h(m)$. Then $a$ has finite support and
$a(n-1)-a(n)=h(n)$, hence $f=s\delta_0+(T-1)a$. Summing these classwise
decompositions gives \eqref{eq:decomp} with finite $A$.

It remains to identify $s$. Within a translation class, the alternating-visit
rule makes the one-period net toggle equal to $+1$ when the stabilised class word
has odd length and $0$ when it has even length. Thus $s\in\{0,1\}$, and $S$ is a
sum of one unit mass for each odd class. The number of odd classes is the net
black growth $\growth=\#\turnR-\#\turnL$, proving the stated form.
\end{proof}

\begin{corollary}\label{cor:sep}
If $\growth=0$, a finite stationary set can be deleted after finitely many periods
to leave a clean finite translator with drift $\drift$.
\end{corollary}

\begin{proof}
Here $S=0$, so $D=(T_\drift-1)A$. Let $b_n$ be the indicator of the black set at
the $n$th period boundary. Translation covariance and telescoping give
\[
b_n=b_0+\sum_{k=0}^{n-1}T_{k\drift}D
   =b_0+(T_{n\drift}-1)A=R+T_{n\drift}A,\qquad R:=b_0-A.
\]
Both $R$ and $A$ are finitely supported, so their supports are disjoint for all
large $n$; because $b_n$ is $\{0,1\}$-valued, each is then the indicator of a
finite set. A nonzero-drift periodic path visits every fixed cell only finitely
often. At a sufficiently late boundary the ant therefore never revisits $R$;
deleting it changes no future turn and leaves the finite pattern $T_{n\drift}A$,
which translates exactly by $\drift$ each period.
\end{proof}

Sections \ref{sec:evenwinding}--\ref{sec:residue} prove that the resulting clean
translator is impossible.

%=======================================================================
\section{The even-winding theorem}\label{sec:evenwinding}
%=======================================================================

\begin{theorem}[\PN{15}: no zero-growth clean translator]\label{thm:evenwinding}
There is no finite black pattern and ant state that evolves in finitely many
steps to exactly one nonzero translate of itself with restored heading and zero
net black growth. Equivalently, no clean finite translator with nonzero drift
exists.
\end{theorem}

\begin{proof}
Suppose $(A,p,q)\mapsto(A+\drift,p+\drift,q)$ in one word of length $N$ with $A$
finite and $\drift\neq0$. By Proposition~\ref{prop:laurent} the drift is
axis-aligned. Since a clean translator has zero net growth,
$N=2\#\turnR$ is even. Each step changes $x+y$ modulo two, so the nonzero drift
coordinate is even. After a quarter-turn rotation write $\drift=(L,0)$ with
$L>0$ even.

\emph{Quotient and the visit count.} Reversibility and translation covariance
give, for all $n\in\Z$ and phases $j$, $A_{nN+j}=A_j+n\drift$ and
$p_{nN+j}=p_j+n\drift$. Quotient the full bi-infinite trace by translation
through $\drift$ to obtain a closed walk on the cylinder
$C_L=(\Z/L\Z)\times\Z$; the checkerboard HV type descends because $L$ is even.
Fix a quotient vertex $v$ with physical representative $z$. Each phase $j$ whose
position lies over $v$ has a unique level $m_j$ with $p_j=z+m_j\drift$, and phase
$j$ in period $-m_j$ visits $z$; conversely every visit to $z$ arises once this
way. Thus the visits to $v$ in one quotient period are exactly the visits to $z$
over the bi-infinite translated orbit. Because at each fixed phase the black
pattern is a finite translate, $z$ is white far in both time directions; its
turns alternate, begin with $\turnR$, end with $\turnL$, and are equally many of
each. Hence every quotient vertex $v$ has
\begin{equation}\label{eq:visitcount}
n_v=2k_v,\qquad k_v\ \turnR\text{'s and }k_v\ \turnL\text{'s},
\end{equation}
with no footprint or gap assumption ($k_v=0$ allowed).

\emph{Edge multiplicities.} For each undirected cylinder edge let its
multiplicity be the number of traversals by the projected one-period trace.
Proposition~\ref{prop:hv} orients each admissible edge uniquely, so at a vertex
one pair of incident edges carries all arrivals and the other all departures:
\begin{equation}\label{eq:balance}
H_{v,-}+H_{v,+}=V_{v,-}+V_{v,+}=2k_v,
\end{equation}
where $H_{v,\pm}$ are the west/east horizontal and $V_{v,\pm}$ the south/north
vertical multiplicities. Reducing the horizontal equality mod $2$ shows
consecutive horizontal edges share parity, so each row $y$ has one horizontal
parity $h_y\in\F_2$. The same propagation up a column, together with finite
vertical support (edges far below/above the trace have multiplicity $0$), forces
\emph{every vertical multiplicity to be even}: $V_e=2G_e$, $G_e\in\Z_{\ge0}$.

\emph{Local parity and the telescope.} A local HV parity computation gives, at
both cell types, $k_v\equiv h_y\pmod2$: at a vertical-arrival cell the horizontal
output has parity $h_y$ and the vertical input is even; at a horizontal-arrival
cell the vertical output is even and the horizontal input has parity $h_y$.
Dividing the vertical part of \eqref{eq:balance} by two gives the exact identity
$k_{x,y}=G_{x,y-1/2}+G_{x,y+1/2}$. Fix a column $x$ and sum over all rows:
\begin{equation}\label{eq:telescope}
\sum_y h_y\equiv\sum_y\bigl(G_{x,y-1/2}+G_{x,y+1/2}\bigr)\equiv0\pmod2,
\end{equation}
because every interior half-edge $G$ occurs twice and the two exterior terms
vanish by finite support.

\emph{Homological contradiction.} Cut the cylinder along a vertical seam between
two adjacent columns. The algebraic winding number $w$ of the projected closed
trace is the signed number of seam crossings; there is one horizontal seam edge
per row, of parity $h_y$, and signs disappear mod $2$, so by \eqref{eq:telescope}
\[
w\equiv\sum_y h_y\equiv0\pmod2 .
\]
But the lift of the closed trace runs from $p$ to $p+(L,0)$ and therefore winds
exactly once, $w=1$. This contradiction proves the theorem for drift $(L,0)$;
quarter-turn rotations, including a half-turn, cover the remaining axis-aligned
drifts.
\end{proof}

The telescoping step \eqref{eq:telescope} is exactly the finite algebraic core
machine-checked in Lean (Appendix~\ref{app:lean},
\texttt{evenWindingParityCore}).

\begin{corollary}[growth is a positive multiple of four]\label{cor:growth4}
Every finite-support periodic highway has net black growth
\begin{equation}\label{eq:growth4}
\growth\in 4\Z_{>0}.
\end{equation}
\end{corollary}

\begin{proof}
Zero drift confines the ant to finitely many positions of one period, excluded by
unboundedness (Theorem~\ref{thm:unbdd}). Negative growth cannot repeat forever
since $b_t\ge0$. A heading-resetting period has $\growth\equiv0\pmod4$
(Proposition~\ref{prop:pathreduction}). By Corollary~\ref{cor:sep} a
$\growth=0$ repeating trace with nonzero drift separates into stationary debris
and a clean moving packet, contradicting Theorem~\ref{thm:evenwinding}. Hence
$\growth>0$ and $\growth\in4\Z_{>0}$.
\end{proof}

The ordinary highway has $\growth=12$, consistent with \eqref{eq:growth4}.

%=======================================================================
\section{The signed mod-four wake-residue theorem}\label{sec:residue}
%=======================================================================

We now give an independent proof of zero-growth exclusion that also produces a
strand-density bound. Recall $\che(x,y)=(-1)^{x+y}$ and the strand bases
$c_1,\dots,c_\growth$ of Theorem~\ref{thm:decomp}.

\begin{theorem}[\PN{22}: signed mod-four wake-residue identity]\label{thm:residue}
Let a finite-support repeating trace have restored heading, nonzero drift
$\drift=(a,b)$, growth $\growth$, and canonical decomposition
\eqref{eq:decomp}. Then:
\begin{enumerate}[label=\textup{(\roman*)}]
\item if $a\neq0$, then for every residue $r\in\Z/|a|\Z$,
\begin{equation}\label{eq:res-a}
\sum_{i:\,c_{i,x}\equiv r\ (|a|)}\che(c_i)\equiv 2\pmod4;
\end{equation}
\item if $b\neq0$, then for every residue $s\in\Z/|b|\Z$,
$\displaystyle\sum_{i:\,c_{i,y}\equiv s\ (|b|)}\che(c_i)\equiv2\pmod4$;
\item if $a=0$ the signed sum over any exact column is $\equiv0\pmod4$, and if
$b=0$ over any exact row.
\end{enumerate}
\end{theorem}

\begin{proof}[Proof]
Fix $a\neq0$ and an arbitrary $\alpha\colon\Z/|a|\Z\to\Z/4\Z$. Put the black-cell
weight
\begin{equation}\label{eq:weight}
w(x,y)=\che(x,y)\,\alpha(x)\in\Z/4\Z .
\end{equation}
Then $4w=0$ and $w$ satisfies the four-corner equation \eqref{eq:cocycle} because
$\che$ changes sign on each unit step; by Proposition~\ref{prop:cocycle} there
are heading potentials $F_q$ with, on an $\turnR$ step
$F_{q+1}(z+u_{q+1})-F_q(z)=-w(z)$ (a black cell of weight $w$ is added) and on an
$\turnL$ step $F_{q-1}(z+u_{q-1})-F_q(z)=+w(z)$.

Composing the north-to-east $\turnR$ equation with the east-to-north $\turnL$
equation gives $F_N(z+(1,1))-F_N(z)=w(z+(1,0))-w(z)=-\che(z)(\alpha(x)+\alpha(x+1))$,
and via $F_S=F_N-2w$ the same increment along $(1,-1)$; since the endpoints differ
by $(0,2)$, $F_N$ is vertically two-periodic. Restored heading gives
$\growth\equiv0\pmod4$, so the period is even and $a+b$ is even; hence translation
by $\drift$ preserves $\che$ and $w(z+\drift)=w(z)$. Telescoping one complete
horizontal residue cycle and using vertical two-periodicity yields the monodromy
\begin{equation}\label{eq:monodromy}
F_q(z+\drift)-F_q(z)\equiv 2\sum_{r\in\Z/|a|\Z}\alpha(r)\pmod4
\end{equation}
for all four headings $q$ and either sign of $a$ (the telescoped total is
$-2\che(z)\sum_r\alpha(r)$, and $-2\equiv2\pmod4$).

Charge conservation over one period reads
\begin{equation}\label{eq:conservation}
\sum_z D(z)\,w(z)+F_q(p+\drift)-F_q(p)=0 .
\end{equation}
Because $w$ is $\drift$-periodic and $(T_\drift A)(z)=A(z-\drift)$, the widget
terms cancel: $\sum_z(T_\drift A)(z)w(z)=\sum_z A(z)w(z)$. Substituting
$D=S+(T_\drift-1)A$ and \eqref{eq:monodromy} into \eqref{eq:conservation} gives
\begin{equation}\label{eq:reduced}
\sum_{i=1}^{\growth}\che(c_i)\,\alpha(c_{i,x})+2\sum_{r}\alpha(r)\equiv0\pmod4 .
\end{equation}
Choosing $\alpha=\mathbf 1_{\{r\}}$ (the indicator of one residue) turns
\eqref{eq:reduced} into $\sum_{i:\,c_{i,x}\equiv r}\che(c_i)+2\equiv0$, i.e.
\eqref{eq:res-a}, using $-2\equiv2\pmod4$. Rotation gives the $b$-statement. For
$a=0$: take a finitely supported exact-column $\alpha\colon\Z\to\Z/4\Z$; then $b$
is even, $w$ is $\drift$-periodic, vertical two-periodicity makes the monodromy
zero, and exact-column indicators give the $\equiv0$ statement; rotation gives
exact rows.
\end{proof}

\begin{corollary}[\PN{16}: strand-density bound]\label{cor:P16}
Every finite-support repeating trace with nonzero drift $\drift=(a,b)$ has
\begin{equation}\label{eq:P16}
\growth\ \ge\ 2\max(|a|,|b|).
\end{equation}
In particular $\growth=0$ is impossible: zero-growth repeating traces do not
exist (an independent proof of Theorem~\ref{thm:evenwinding}). If a residue fiber
contains exactly two strands, both bases share checkerboard parity.
\end{corollary}

\begin{proof}
A sum of $m$ signs $\pm1$ that is $\equiv2\pmod4$ is nonzero and has $m$ positive
and even. Hence each of the $|a|$ residue classes mod $|a|$ (and each of the $|b|$
classes mod $|b|$) contains at least two strand bases, giving $\growth\ge2|a|$ and
$\growth\ge2|b|$. If $\growth=0$ some class is empty while $\drift\neq0$,
contradicting \eqref{eq:res-a}. A two-strand fiber has $\che$-sum $\pm2$, forcing
equal signs.
\end{proof}

Unlike the even-winding proof, Theorem~\ref{thm:residue} needs neither the axis
alignment of Proposition~\ref{prop:laurent} nor the clean-packet reduction; it
applies directly to every repeating trace.

%=======================================================================
\section{Collision-chain parity}\label{sec:interlace}
%=======================================================================

Encode a finite cell set $D$ as edges of the bipartite row--column incidence
graph: cell $(x,y)$ is the edge $C_xR_y$. At a time $s$, assign the token
$\eta_s=C_{p_{s,x}}$ if $q_s$ is vertical and $\eta_s=R_{p_{s,y}}$ if $q_s$ is
horizontal.

\begin{theorem}[Two-endpoint collision-chain parity]\label{thm:collision}
Over $\F_2$, the boundary (set of odd-degree vertices) of the row--column incidence
graph of $B_s\triangle B_t$ equals $\eta_s+\eta_t$. Intermediate tokens telescope
across arbitrary collision checkpoints; hence at most one defect component is
non-Eulerian, and a same-pose return ($\eta_s=\eta_t$) has only Eulerian defect
components.
\end{theorem}

\begin{proof}
One step toggles the single edge $C_{p_{s,x}}R_{p_{s,y}}$, changing the degree
parity of exactly its two endpoints. If the pre-step heading is vertical then
$\eta_s=C_{p_{s,x}}$ and the post-step heading is horizontal with unchanged row,
so $\eta_{s+1}=R_{p_{s,y}}$; for a horizontal pre-step heading the roles are
reversed. Hence $\eta_s+\eta_{s+1}=C_{p_{s,x}}+R_{p_{s,y}}$, exactly the one-step
boundary contribution. Summing over a run telescopes to $\eta_s+\eta_t$. By the
handshaking lemma, every connected component has an even number of odd-degree
vertices. If $\eta_s\ne\eta_t$, the graph has exactly those two odd vertices,
so they lie in the same component, which is the unique non-Eulerian component. If
$\eta_s=\eta_t$, the boundary is empty and every component is Eulerian.
\end{proof}

This identity gives a four-cell lower bound on completed lifetimes. Let $s$ be
immediately after an $\turnR$ visit to $z=(x,y)$ and $t$ immediately before its
next ($\turnL$) visit. Then $z$ is black at both times, so $C_xR_y$ is absent from
$G:=B_s\triangle B_t$. If arrivals at $z$ are vertical, the $\turnR$ step leaves
$z$ horizontally along row $y$, so $\eta_s=R_y$, while $\eta_t=C_x$; for
horizontal arrivals the roles reverse. Hence $\partial G=C_x+R_y$. The two
vertices are joined in $G$ by a path avoiding $C_xR_y$; in the simple bipartite
incidence graph such a path has at least three edges. Adjoining $C_xR_y$ closes a
cycle of length at least four, so the completed lifetime involves at least four
distinct cells.

Circuit-partition and interlacement formalisms may help compress the remaining
chronological constraints, but no touch-graph rank theorem is claimed here.

%=======================================================================
\section{Transverse rigidity and all-period narrow-width exclusions}\label{sec:transverse}
%=======================================================================

The constraints above are incidence statements. Section~\ref{sec:gap} records
countermodels that satisfy every one of them and are still not realisable, so no
sharpening of the residue equations alone can close the classification problem. This
section proves statements of a different kind, using the \emph{order} in which a
cell is visited and exact local dynamics rather than only aggregate counts.  The
width-two and computer-assisted width-four exclusions both hold at every period.

Assume the drift is diagonal: $\drift=(a,b)$ with $|a|=|b|=m\ge1$. A quarter-turn
rotation sends $\drift$ to $(-b,a)$, hence $ab$ to $-ab$, and commutes with the ant's
rule; since highways are identified up to quarter-turn rotation we normalise to
$ab<0$ and set $t(x,y)=x+y$, so $\drift=me$ with $e=\pm(1,-1)$. Then $t(\drift)=0$,
so $t$ is constant on translation classes, and every unit step changes $t$ by exactly
$\pm1$: the increments are $+1$ for east and north, $-1$ for west and south. (For an
axis drift the transverse functional is $x$ or $y$, which two moves fix, and
everything below fails.) Put $T=\max_i t(p_i)$, $T'=\min_i t(p_i)$ and $W=T-T'$. All
cells of a level line $t=c$ share the checkerboard parity of $c$, so by the HV
partition each level line is entirely horizontal-arrival or entirely vertical-arrival.

The normalisation is essential rather than cosmetic. In the frame $ab>0$ one must take
$t=x-y$, whose $t$-increasing headings are east and \emph{south}; there the argument
below eliminates the horizontal case and the extremal lines are \emph{vertical}. Only
the axis-free conclusions---every extremal turn is $\turnR$, extremal cells are entered
at most once, $W$ is even---read identically in both frames.

\begin{theorem}[extremal-line rigidity]\label{thm:extremal}
For a finite-support periodic highway with diagonal drift, normalised as above, both
extremal level lines $t=T$ and $t=T'$ are horizontal; the arrival is east on $t=T$
and west on $t=T'$; every turn on them is $\turnR$; every extremal cell reached by
the periodic tail is entered exactly once during that tail, so every translation
class meeting an extremal line is a singleton and its tail cells are left black and
never revisited thereafter; and $W$ is even.
\end{theorem}

\begin{proof}
Let $I$ be a class meeting $t=T$ and let $z_n=b_I+n\drift$, all on that line since
$t(\drift)=0$. Take $n$ with $z_n$ reached by every phase of $I$, so the turn sequence
at $z_n$ is the stabilised word of $I$ and every visit has a predecessor. Arrivals at
$z_n$ increase $t$ and departures do not. Recall $\turnR$ sends north, east, south,
west to east, south, west, north, and $\turnL$ reverses this. Were the line vertical,
the arrival would be north, whose clockwise image east also increases $t$ and is
forbidden, while its anticlockwise image west is allowed; every turn would be
$\turnL$, so the stabilised word would begin with $\turnL$, contradicting the
criterion. So the line is horizontal with arrival east, whose anticlockwise image
north is forbidden and clockwise image south allowed: every turn is $\turnR$. At
$t=T'$ arrivals decrease $t$ and departures must increase it; the vertical case gives
arrival south, clockwise image west forbidden, anticlockwise east allowed, hence all
$\turnL$ and the same contradiction, so that line is horizontal with arrival west and
every turn $\turnR$. In both cases the stabilised word is a block of $\turnR$'s and
must alternate, so $|I|=1$. Both extremal lines being horizontal, they share a
checkerboard parity, so $T\equiv T'\pmod2$.
\end{proof}

\begin{theorem}[no width-two highway]\label{thm:width2}
No finite-support periodic highway with diagonal drift has $W=2$; hence $W$ is even
and $W\ge4$.
\end{theorem}

\begin{proof}
Label the three lines $0,1,2$ for $t=T,T-1,T-2$; lines $0,2$ are extremal and
horizontal, line $1$ is vertical. Each step changes $t$ by $\pm1$, so the ant
alternates between line $1$ and $\{0,2\}$. Orienting the longitudinal axis so that
$e=(1,-1)$, the excursions are forced by Theorem~\ref{thm:extremal}: from a line-$1$
cell $z$ the upward excursion enters line $0$ heading east and leaves heading south,
carrying $z\mapsto z+e$; the downward one enters line $2$ heading west and leaves
heading north, carrying $z\mapsto z-e$. Index line-$1$ cells by $\ell\in\Z$ along $e$.

An upward excursion out of $\ell$ consumes the line-$0$ cell $U_\ell$ and a downward
one the line-$2$ cell $D_\ell$; distinct sites consume distinct cells, each entered at
most once by Theorem~\ref{thm:extremal}. The gap $(\ell,\ell+1)$ is crossed upward
only from $\ell$ and downward only from $\ell+1$, so it is crossed at most once in
each direction. Choose a cyclic phase at a visit to line $1$, and let $s_0$ be that
starting site. Drift $+m$ sends the site to
$+\infty$, so every gap above $s_0$ has net crossing number exactly $+1$; with each
count at most one this forces one upward and no downward crossing. Hence every site
$\ell>s_0$ is entered exactly once, from $\ell-1$, and departs to $\ell+1$.

Entering from $\ell-1$ completes an upward excursion, which leaves line $0$ heading
south, so the ant arrives heading south; departing to $\ell+1$ begins an upward
excursion, so it leaves heading east. The anticlockwise image of south is east, so the
unique turn at $\ell$ is $\turnL$. The class of such a site has cells $\ell+nm$, almost
all above $s_0$ and hence entered once, so it is a singleton whose stabilised word is
the single letter $\turnL$---contradicting the criterion. \end{proof}

We record the two errors this argument passed through, since both are natural traps.
Writing the transverse functional with a case split and then arguing in one frame gives
a false ``both extremal lines are horizontal'' in the rotated frame, where they are
vertical; and asserting that a site's first visit turns $\turnR$ is wrong, because the
criterion constrains only stabilised words and the seed construction blackens cells
below their class's maximal level precisely so that some first turns are $\turnL$. The
proof above avoids the first by normalising and using explicit compass directions, and
the second by counting edge crossings rather than site visits.

The residue theorem and extremal parity also give a useful general inequality.  If
the normalised drift is $(m,-m)$, growth is $g$, and transverse width is $W$, then
\begin{equation}\label{eq:widthgrowth}
2m\le g\le mW.
\end{equation}
Indeed a translation class is labelled by $(t,x\bmod m)$.  In each fixed
$x$-residue the signed residue identity makes the number of odd classes positive and
even.  There is at most one on each of the $W+1$ levels; since $W$ is even, the
largest possible even number is $W$.  Summing over the $m$ residues proves the upper
bound, and the strand-density theorem gives the lower bound.  For the standard
highway $(m,g,W)=(2,12,10)$ this yields only $W\ge6$.

The width-four case can nevertheless be excluded directly.  Rotate by $180^\circ$
if needed, take drift $(m,-m)$ with $m>0$, shift the transverse range to
$t=x+y\in\{0,1,2,3,4\}$, and view each fixed-$x$ column as a five-bit tape symbol.
At a half-integer boundary record, in chronological order, the transverse line of
the destination cell after each horizontal crossing.  Directions alternate and are
implicit: rightward entries are $2$ or $4$, leftward entries are $0$ or $2$.
Extremal rigidity permits rightward $4$ and leftward $0$ at most once on each
boundary.

\begin{lemma}[exact blank-column graph]\label{lem:blankcolumn}
For one initially white five-cell column, among boundary sequences that alternate
right/left beginning and ending right, use only rightward labels in $\{2,4\}$ and
leftward labels in $\{0,2\}$, and use rightward $4$ and leftward $0$ at most once on
each boundary, the complete list of compatible left-to-right pairs is
\[
\begin{array}{c|l}
202&4\\
4220222&2,\ 202,\ 20222,\ 2022224\\
2042222&2,\ 202,\ 20222,\ 2022224\\
202242222&422,\ 42202,\ 4220222.
\end{array}
\]
This twelve-edge graph is acyclic and has longest directed path three.
\end{lemma}

\begin{proof}[Exact computer-assisted verification]
Literal simulation tracks the $32$ colour masks, current entry line and heading,
and four flags for the one allowed extreme event on either side.  The complete
recursion reaches $22$ entry states, has no reachable state cycle, and terminates by
depth eight.  A second implementation, using compass headings and sets of black rows
and sharing no classifier code, gives the same twelve pairs.  It also extracts two
far cuts from the documented standard-highway seed and exactly replays the initially
white intervening column.  Assign ranks $3,2,1,0$ respectively to
$202242222$, the pair $4220222,2042222$, the vertex $202$, and all other displayed
vertices.  Every edge lowers rank; this last finite implication is checked in Lean.
The Lean file does not formalise completeness of the two Python enumerations.
\end{proof}

\begin{theorem}[no width-four highway]\label{thm:width4}
No finite-support periodic highway with diagonal drift has transverse width $W=4$.
Hence every such highway has even width $W\ge6$.
\end{theorem}

\begin{proof}
At the beginning of the periodic tail only finitely many cells have been visited and
only finitely many cells are black, by finite initial support.  Hence all sufficiently
far columns are untouched and white.  At any far boundary the
longitudinal coordinate is eventually $nm$ plus a bounded phase offset; its complete
crossing sequence is therefore finite, nonempty, and alternates right, left,
$\ldots$, right.  For every far column, its left and right sequences form an edge of
Lemma~\ref{lem:blankcolumn}.  Infinitely many consecutive white columns would give an
infinite directed path, contradicting the rank bound.
\end{proof}

There is an independent arithmetic view of the width-four geometry.

\begin{proposition}[width-four masks and period bound]\label{prop:width4arithmetic}
Let a finite-support periodic highway have transverse width four and normalised drift
$(m,-m)$, $m>0$.  For each longitudinal residue, the set of levels carrying odd translation classes
is one of
\[
\{0,2\},\ \{1,3\},\ \{0,4\},\ \{2,4\},\
\{0,1,2,4\},\ \{0,2,3,4\};
\]
and for each diagonal coordinate residue the corresponding set belongs to the same
six-element list.  If $h$ of the $m$
longitudinal residues use a four-element mask, then $g=2m+2h$ and
$h\equiv m\pmod2$.  If $P$ is the macro period and $N=2P$ the ant period, then
\[
P\le7m,\qquad N\le14m,\qquad P\equiv m\pmod2.
\]
\end{proposition}

\begin{proof}
Let $o_{j,r}$ indicate that the class on level $j$ and longitudinal residue $r$ is
odd.  The signed residue identity is
$o_{0,r}-o_{1,r}+o_{2,r}-o_{3,r}+o_{4,r}=2\pmod4$.
Exhausting five binary values gives exactly the six displayed masks; applying the
other coordinate residue identity gives the diagonal version.  Summing mask sizes
gives $g=2m+2h$, and $g\in4\Z$ gives the parity of $h$.

For the period bound, cyclically start immediately after a bottom-extremal visit and
observe the ant only on the odd levels.  The forced bottom turn puts it in state
$LN$, and one translated period returns to $LN$; width four ensures that a top-line
transition also occurs.  Write the four states
$LN,LS,UN,US$ for lower/upper level and north/south arrival.  The twelve possible
two-turn macro transitions are
\[
\begin{array}{c|ccc}
LN&(LS,+1,RR)&(UN,+1,RL)&(LN,-1,LR)\\
LS&(LN,-1,RR)&(LS,+1,LR)&(UN,+1,LL)\\
UN&(US,+1,RR)&(UN,-1,LR)&(LS,-1,LL)\\
US&(UN,-1,RR)&(LS,-1,RL)&(US,+1,LR).
\end{array}
\]
Let $x$ count $LN\to LN$, let $y$ count $US\to US$, let $c$ count $LS\to LN$,
and let $f$ count $UN\to US$.  If $a,b$ count $LN\to LS,LN\to UN$ and $i,j$
count $US\to UN,US\to LS$, flow conservation gives $c=a+b$ and $f=i+j$.
Extremal singleton rigidity gives $x+c\le m$ and $f+y\le m$.  If $L_1,L_3$ are
the lower and upper odd-level left-turn totals, the P3 condition that each
stabilised class alternates and begins with $R$ makes the total $R-L$ excess on
each odd level nonnegative.  The lower right-turn total is $a+b+c=2c$ and the
upper one is $f+i+j=2f$, so $L_1\le2c$ and $L_3\le2f$.  Flow conservation and
the displacement read from the table give
\[
P=2c+2f+L_1+L_3,\qquad m=L_1-L_3-2x+2y.
\]
Thus $L_1+L_3\le3m-6x+2y$, and substitution yields
$P\le7m-8x\le7m$.  Every macro move changes the longitudinal coordinate by
$\pm1$, so $P\equiv m\pmod2$.
\end{proof}

Combined exact searches cover the entire theorem-bounded range for $1\le m\le9$:
the fixed-drift engine supplies the
full range for $m\le6$ and the upper ranges for $m=7,8,9$, while the general
period-$\le48$ and unrestricted width-four searches supply their lower ranges.  The
result is zero hits throughout.  The $m=10$ run is
partial through $P=66$, and its unfinished rows are not counted.  These arithmetic
and search results are cross-checks and structural data; Theorem~\ref{thm:width4}
does not depend on them or on the residue theorem.

The proof has a sharp width boundary.  At widths six and eight the local recursion
contains a reachable two-state oscillator: black rows $\{1,2\}$ with an eastward
entry on row $2$ change to $\{1,3\}$ and exit eastward, and a westward return on row
$2$ restores the first state.  The reused cells alternate their turns and neither
extreme repeats.  This is not a highway, but it shows that one-column acyclicity does
not extend unchanged.  Nor does a finite column alphabet alone make an unbounded
tape finite-state; wider strips require a genuine bounded-memory theorem.

%=======================================================================
\section{Exact computational theorems}\label{sec:computation}
%=======================================================================

All computations use exact integer coordinates, exact finite sets, and no
probabilistic equality tests. The residue-free Java variant certifies the
period-$34$--$48$ exclusion. Its strand- and residue-pruned variants share one
enumeration framework and are consistency checks, not independent
implementations. A separately written Python verifier independently reproduces
the complete exclusion through length $18$ and audits the exact criterion.

\begin{theorem}[gateway seed]\label{thm:gateway}
The blank orbit first enters its repeating trace at step $9977$; the trace has
least period $104$ and normalised drift $(-2,2)$. Applying
Theorem~\ref{thm:criterion} to this word yields the $13$-cell black seed
\[
\begin{aligned}
&(-6,0),(-5,0),(-5,1),(-4,0),(-4,2),(-3,-1),(-3,2),\\
&(-2,-1),(-2,1),(-1,1),(0,0),(0,1),(0,2),
\end{aligned}
\]
which, with a north-facing ant at the origin, produces the $104$-step highway
immediately; the semi-infinite leading corridor is white and the gateway
induction recurs translated by the drift.
\end{theorem}

\begin{theorem}[bounded-period exclusion]\label{thm:periodexcl}
No finite-support periodic highway of nonzero drift has period $\le 48$. The
verification proceeds by:
\begin{enumerate}[label=\textup{(\alph*)}]
\item a complete first-$\turnR$, cyclic-phase-normalised enumeration through
period $32$ ($46{,}185{,}421$ nodes at depths $6$--$32$ over $32$ prefix shards,
plus a $59$-node unprefixed baseline through depth $5$; $46{,}185{,}480$ nodes in
total). Every positive-growth heading-resetting word contains $\turnR$, so a cyclic
shift places it in this enumeration; shifting changes only the phase on the same
periodic orbit and preserves finite-support realisability;
\item rank-by-rank audited positive-growth searches at periods $34,\dots,48$ with the
\emph{residue-free variant}, which uses no consequence of Theorem~\ref{thm:residue} and
applies the exact criterion to \emph{every} positive-growth nonzero-drift leaf, using
Corollary~\ref{cor:growth4} to skip zero growth; odd periods restore no heading.
The exact node and leaf counts are:
\begin{center}
\small
\begin{tabular}{@{}rrr@{}}
\toprule
Period & Nodes & Leaves\\
\midrule
$34$ & $40{,}418{,}033$      & $6{,}989{,}418$\\
$36$ & $108{,}995{,}287$     & $18{,}928{,}333$\\
$38$ & $293{,}763{,}689$     & $51{,}386{,}767$\\
$40$ & $791{,}801{,}484$     & $139{,}000{,}900$\\
$42$ & $2{,}133{,}302{,}151$ & $376{,}781{,}423$\\
$44$ & $5{,}748{,}138{,}964$ & $1{,}018{,}293{,}035$\\
$46$ & $15{,}483{,}269{,}352$& $2{,}757{,}152{,}898$\\
$48$ & $41{,}710{,}394{,}384$& $7{,}446{,}719{,}550$\\
\bottomrule
\end{tabular}
\end{center}
At period $48$ all $2^{16}$ prefix ranks were covered with no node cap. The
exact criterion was evaluated at every leaf and returned zero certificates at
every period; every rank was covered exactly once.
Strand-pruned and residue-pruned variants, which do assume
Theorem~\ref{thm:residue}/\ref{cor:P16}, return zero certificates on the strictly
smaller trees they explore and serve as cross-checks only.
\end{enumerate}
Together with Corollary~\ref{cor:growth4} (which excludes zero growth at every
period) this excludes every finite-support periodic highway of period at most
$48$.
\end{theorem}

These theorems are exact for the stated finite ranges only. They bound the least
period of a possible nonstandard highway, whereas Conjecture~\ref{conj:main}
quantifies over every finite configuration and every period; this finite search
does not exhaust that range.

%=======================================================================
\section{Remaining gaps and open problems}\label{sec:gap}
%=======================================================================

The results above leave two logically separate problems. The first is a
universal entrance problem; the following are two sufficient formulations,
which are not claimed to be equivalent.

\begin{problem}[bounded-retreat lemma]\label{prob:retreat}
For every finite-support initial configuration there exist a diagonal direction,
a time $T$, and a width $W<\infty$ such that after $T$ the ant never retreats
more than $W$ layers behind its current record frontier and all transverse
excursions stay in a width-$W$ corridor.
\end{problem}

\begin{problem}[bounded active core]\label{prob:core}
For every finite-support initial configuration there are a finite window
$W\subset\Z^2$, a time $T$, and translations $v_t$ such that, for every $t\ge T$,
the ant and every cell or Tait-boundary component capable of influencing the future
orbit lie in $v_t+W$; all components outside that translated window are permanently
inert and never reabsorbed.
\end{problem}

Under either formulation, the colours in a bounded frontier window, the ant's
position and heading form a finite state modulo translation, so determinism
forces eventual periodicity. This would not identify the resulting cycle with
the standard highway; that is the separate periodic-classification problem
below. No proof of either entrance formulation is known, and the following
obstructions show why:

\begin{itemize}[leftmargin=1.4em]
\item Unboundedness (Theorem~\ref{thm:unbdd}) and fixed-domain escape do not
imply bounded retreat: an orbit may keep enlarging a two-dimensional region while
escaping every fixed box.
\item No nonconstant affine combination of $(E,V,k,\beta)$---in particular no single
scalar Tait invariant among them---is monotone (Proposition~\ref{rem:nomonotone}).
\item A local pumping argument is invalid: equal local cores can be joined
through different distant connectivity, changing a later merge into a split
(Theorem~\ref{thm:conjugacy}). A valid finite-state proof must first establish
permanent separation from the discarded graph or a finite connectivity summary.
\item At fixed growth $\growth=4$ the signed-residue skeletons of
Theorem~\ref{thm:residue} leave unbounded free parameters; six exact
heading-reset countermodels satisfy every incidence/residue condition yet fail
the cross-period criterion, so these displayed incidence and residue conditions
alone cannot finish the periodic branch.
\end{itemize}

\begin{problem}[periodic classification]\label{prob:periodic}
Prove that every positive-growth periodic highway admitted by finite support agrees
in path and turn trace with the standard period-$104$ highway up to lattice
translation, quarter-turn rotation, and cyclic phase shift.
One identified route requires a single-tour chronology theorem valid for unbounded period; the
bounded-period exclusions of Section~\ref{sec:computation} are only a finite
frontier.  Theorem~\ref{thm:width4} removes one unbounded family, but widths six and
above, non-diagonal translators, and uniqueness within the surviving cases remain.
\end{problem}

A genuine \emph{disproof} would output an infinite certificate: either a
nonstandard periodic word passing Theorem~\ref{thm:criterion} (disproving
uniqueness of the highway), or a finite self-enlarging Tait gadget whose state
recurs with a parameter increased (an induction proof of nonperiodic growth). No
such certificate has been found.

%=======================================================================
\section{Summary}
%=======================================================================

We have assembled a coherent theory of the finite-support Langton ant: an exact
decision criterion for periodic highways, a Tait-graph conjugacy with surgery
calculus, two independent proofs that no zero-growth clean translator exists
(hence every highway grows a wake of size $\in4\Z_{>0}$), the strand-density
bound $\growth\ge2\|\drift\|_\infty$, exhaustive exclusion of all periodic
highways of period $\le48$, extremal-line rigidity, and all-period exclusions of
diagonal transverse widths two and four.  The width-four theorem is
computer-assisted, with two independent local enumerators and a Lean-checked graph
rank certificate.  We also give reductions relevant to two remaining problems:
universal entrance and all-period classification. Selected algebraic kernels are
machine-checked in Lean~4; the geometric extraction is not formalized end to end
(Appendix~\ref{app:lean}).
Conjecture~\ref{conj:main} remains open; neither remaining problem is solved
here.

\appendix
%=======================================================================
\section{Lean 4 formalization of selected algebraic kernels}\label{app:lean}
%=======================================================================

A dependency-light Lean~4 project (pinned to Lean $4.32.0$, standard library
only) machine-checks the exact transition kernel and the algebraic invariants
used above. All audited theorems depend only on the standard axioms
\texttt{propext}, \texttt{Classical.choice}, \texttt{Quot.sound}; there are no
\texttt{sorry}s, no custom axioms, and no \texttt{native\_decide}. The checked
statements include:
\begin{itemize}[leftmargin=1.4em]
\item the finite-support transition kernel and its toggle lemmas
(the departed cell flips; every other cell is unchanged);
\item the heading/black-count reset theorem
(Proposition~\ref{prop:pathreduction}, $\bmod 4$);
\item the finite parity telescope \texttt{evenWindingParityCore}: adjacent
half-flow parities bookended by zero exterior values have zero aggregate seam
parity---the core \eqref{eq:telescope} of Theorem~\ref{thm:evenwinding}---and the
contradiction wrapper \texttt{noOddWindingFromFiniteHalfFlow};
\item the four-corner admissibility and mod-four fiber arithmetic of
Theorem~\ref{thm:residue}: $E_r+2=0\Rightarrow E_r=2$ in $\Z/4$, a residue
fiber $=2$ has a positive even number of bases, residue fibers partition the base
list, and one identity per residue implies $\growth\ge2n$;
\item finite potential telescoping, checker-cycle monodromy $2\sum\alpha$, and
translated-widget cancellation (the steps of \eqref{eq:monodromy}--\eqref{eq:reduced});
\item collision-chain endpoint parity (Theorem~\ref{thm:collision}) and the
separately formalized exact directed-pose row/column $\Z/4$ discrepancy charges.
\item the width-four arithmetic kernels (six-mask enumeration, period--drift
inequality, period parity, and the primitive arithmetic reduction), and the finite
rank certificate for Lemma~\ref{lem:blankcolumn}.  Lean checks every listed graph
edge, not completeness of the independently generated edge table.
\end{itemize}
The formalization deliberately isolates, as an explicit interface, the two
inputs it does not derive from the kernel: the endpoint normal form and the
lifted potential geometry of a genuine periodic ant trace certificate. Consequently neither
Theorem~\ref{thm:residue}/\ref{cor:P16} nor Conjecture~\ref{conj:main} is claimed
as an unconditional Lean theorem; what is certified is the algebraic layer
conditional on a structured trace certificate, together with the growth bound it
yields.

\section*{Reproducibility}
The exact search engines, the periodic-trace checker, the Lean project, and all
period-exclusion records with SHA-256 manifests accompany this work in the
supplementary repository \cite{Repo}. The standard word,
the $13$-cell seed, and every period-$\le48$ shard are regenerable from the
sources. The period-$34$--$48$ certifying shards record
\texttt{search\_complete} and \texttt{node\_cap=null}. The period-$\le32$
prefix shards record a $10^8$ safety cap, which every shard left unreached (the
largest used $2{,}623{,}020$ nodes). Released audit scripts verify rank coverage,
aggregate counters, and zero-hit outcomes; agreement among the three Java
variants is a consistency check rather than independent replication.
The width-four supplement additionally contains the two independent blank-column
classifiers, their twelve-edge JSON certificate, the Lean rank module, the complete
fixed-drift summary for $m\le9$, and the explicitly partial $m=10$ record.

%=======================================================================
%  Bibliography embedded so this file compiles stand-alone in one paste
%  (no external references.bib and no BibTeX step required).
%=======================================================================
\begin{thebibliography}{99}

\bibitem{BunimovichTroubetzkoy1992}
L.~A. Bunimovich and S.~E. Troubetzkoy.
\newblock Recurrence properties of {L}orentz lattice gas cellular automata.
\newblock \emph{J. Stat. Phys.} \textbf{67} (1992), no.~1--2, 289--302.

\bibitem{Etse2026confined}
K.~R. Etse.
\newblock How long can the escaping ant be confined?
\newblock arXiv:2606.26677 (2026).

\bibitem{Gajardo2014symbolic}
A.~Gajardo.
\newblock A symbolic projection of {L}angton's ant.
\newblock \emph{Discrete Math. Theor. Comput. Sci. Proc.} AB (DMCS'03)
  (2003), 57--68.
\newblock doi:10.46298/dmtcs.2312.

\bibitem{GajardoLutfallaRao2024}
A.~Gajardo, V.~H. Lutfalla, and M.~Rao.
\newblock Ants on the highway.
\newblock arXiv:2409.10124 (2024); \emph{Natural Computing} \textbf{24} (2025), 497--509.

\bibitem{GajardoMoreiraGoles2002}
A.~Gajardo, A.~Moreira, and E.~Goles.
\newblock Complexity of {L}angton's ant.
\newblock \emph{Discrete Appl. Math.} \textbf{117} (2002), no.~1--3, 41--50.
\newblock arXiv:nlin/0306022.

\bibitem{Lutfalla2025sideways}
V.~H. Lutfalla.
\newblock Sideways on the highways.
\newblock arXiv:2505.05426 (2025).

\bibitem{Repo}
A.~Jillhewar.
\newblock Langton's ant: structural constraints, exact searches, and a Lean
formalization (supplementary code, data, and Lean project).
\newblock Zenodo, 2026. Concept \textsc{doi}:\,\href{https://doi.org/10.5281/zenodo.21381637}{10.5281/zenodo.21381637}.
\newblock Archived from \url{https://github.com/Atharva12456/langtons-ant-highway}.
\newblock The concept DOI resolves to the latest archived release; a version DOI should
be used to pin an exact release, together with its file-level SHA-256 manifest.

\end{thebibliography}

\end{document}
